\documentclass[11pt,letterpaper]{article}

%==================================================
% PACKAGES
%==================================================

\usepackage[
    margin=2cm
]{geometry}

% Make the PDF readable by automated recruiting systems
\input{glyphtounicode}
\pdfgentounicode=1
\usepackage[T1]{fontenc}

\usepackage[dvipsnames]{xcolor}
\usepackage[
    pdftitle={CV - Antonio Rubio Castaneda},
    pdfauthor={Antonio Rubio Castaneda},
    pdflang={es-MX},
    hidelinks
]{hyperref}
\usepackage{titlesec}
\usepackage{enumitem}

% Sourcesanspro ships with almost every TeX distribution
\usepackage[default]{sourcesanspro}
\renewcommand{\familydefault}{\sfdefault}

% No hyphenation anywhere. The slack that hyphens used to absorb goes into the
% right margin instead, so nothing overflows.
\usepackage[none]{hyphenat}
\usepackage{microtype}

\usepackage{lastpage}

%==================================================
% SETTINGS
%==================================================

\pagestyle{plain}

% Ragged right edge. With hyphenation off this is what keeps long words from
% being pushed past the right margin by justification.
\raggedright

\setlength{\parindent}{0pt}
\setlength{\parskip}{4pt}

\setlist[itemize]{
    leftmargin=18pt,
    itemsep=2pt,
    topsep=4pt
}

%==================================================
% HEADINGS
%==================================================

\titleformat{\section}
{\Large\bfseries}
{}
{0em}
{}
[\titlerule]

\titlespacing{\section}
{0pt}
{14pt}
{8pt}

%==================================================
% COMMANDS
%==================================================

% Entry header: name on the left, date or link on the right. Used for both jobs
% and projects, they render the same.
% \makebox is required: under \raggedright the \rightskip glue competes with
% \hfill and the date would not reach the right margin.
% \addvspace instead of \vspace so the section heading's own space is not
% doubled on the first entry, it takes the max of the two.
% The % at line ends kills the spurious spaces the macro used to emit
\newcommand{\entry}[2]{%
\addvspace{8pt}%
\makebox[\linewidth]{\textbf{#1}\hfill{\small#2}}%
}

\newcommand{\position}[1]{%
{\itshape #1}%
}

%==================================================
% DOCUMENT
%==================================================

\begin{document}

\makeatletter
\renewcommand{\@oddfoot}{\hfil \thepage\ de \pageref{LastPage}\hfil}
\makeatother

%==================================================
% HEADER
%==================================================

\begin{center}

{\fontsize{28}{30}\selectfont\bfseries
Antonio Rubio Castañeda}

\vspace{12pt}

\href{mailto:admin@antruc.dev}{admin@antruc.dev}%
\hspace{1cm}%
\href{https://antruc.dev}{antruc.dev}%
\hspace{1cm}%
\href{https://github.com/antruc}{github.com/antruc}

\end{center}

%==================================================
% PROFILE
%==================================================

\section{Perfil Profesional}

Desarrollador de software autodidacta con más de ocho años de experiencia en
desarrollo web y aplicaciones Android. Actualmente dirijo el área de informática
de CONSAEFA, donde me dedico a automatizar los procesos internos, administrar la
infraestructura tecnológica de la empresa y desarrollar software.

%==================================================
% EXPERIENCE
%==================================================

\section{Experiencia Profesional}

\entry{CONSAEFA S.C.}{Junio 2024 -- Actualidad}

\position{Responsable de Informática}

\begin{itemize}
    \item Construí un sistema de inventario con altas, bajas y consultas, primero en VBA y migrado después a aplicación de Android y de escritorio.
    \item Desarrollé aplicaciones de registro en campo que entregan los datos ya consolidados y listos para su análisis, sin transcripción posterior.
    \item Mejoré el análisis SIG con una librería de Python propia y, sobre ella, el proceso completo de clasificación de superficies de predios, de las curvas de nivel a la capa final.
    \item Automaticé el cálculo y los reportes de gastos semanales, antes hechos a mano.
    \item Mantengo la presencia digital de la empresa, que abarca la intranet en SharePoint, el sitio web corporativo y las redes sociales.
    \item Doy soporte técnico a los equipos de cómputo y al equipo de oficina.
\end{itemize}

\entry{Desarrollador Web Freelance}{Enero 2022 -- Junio 2024}

\position{Desarrollo Web y Sistemas Embebidos}

\begin{itemize}
    \item Construí sitios web para negocios locales con WordPress y tecnologías web modernas, desde el diseño hasta la publicación y el mantenimiento.
    \item Desarrollé aplicaciones web a medida con JavaScript según los requerimientos de cada cliente.
    \item Fabriqué dispositivos de propósito específico con Arduino y Raspberry Pi Pico.
    \item Produje y edité contenido gráfico y audiovisual para clientes.
\end{itemize}

\entry{BudaCat}{Septiembre 2018 -- Diciembre 2021}

\position{Desarrollador de Aplicaciones Android}

\begin{itemize}
    \item Desarrollé y publiqué aplicaciones Android que suman más de 4 millones de descargas y una calificación promedio de 4.5 estrellas.
    \item Diseñé la interfaz, la identidad visual y los recursos gráficos de los productos, y gestioné sus campañas de adquisición en Google Ads.
\end{itemize}

\entry{Comex}{Agosto 2014 -- Abril 2016}

\position{Encargado de Tienda}

\begin{itemize}
    \item Administré la operación diaria de la sucursal, desde la atención a clientes hasta el inventario, la caja y la facturación electrónica.
\end{itemize}

%==================================================
% PROJECTS
%==================================================

\section{Proyectos}

\entry{cfasig}{\href{https://pypi.org/project/cfasig/}{pypi.org/project/cfasig}}

Librería de Python sobre GeoPandas, Shapely y Rasterio que lleva el análisis SIG
de las ventanas de diálogo de ArcGIS a scripts reproducibles, con equivalencia
directa a las operaciones de ArcToolbox.

\entry{Depurador de Cámara Trampa}{\href{https://github.com/antruc/depuradordecamaratrampa}{github.com/antruc/depuradordecamaratrampa}}

Aplicación de escritorio que revisa lotes de imágenes de cámaras trampa y
conserva solo aquellas donde aparece fauna, con Ultralytics YOLO sobre ONNX
Runtime. Procesamiento en flujo para lotes grandes sin cargar todo en memoria.
Python, Toga y Briefcase; empaquetada para Linux y Windows.

\entry{Bongo Cat}{\href{https://github.com/antruc/bongo-cat}{github.com/antruc/bongo-cat}}

Instrumento musical táctil para Android, publicado en 2018 y en mantenimiento
activo desde entonces: más de un millón de descargas, 10\,000 reseñas y 4.1
estrellas en Google Play. Dieciocho instrumentos con carga y reproducción de
audio de baja latencia sobre Web Audio, para que el sonido responda al toque sin
retraso perceptible. Aplicación web empaquetada con Cordova y compilada con
esbuild; migración a Capacitor en curso.

\entry{tumba}{\href{https://www.npmjs.com/package/tumba}{npmjs.com/package/tumba}}

Herramienta de línea de comandos para cifrar y descifrar archivos con AES-GCM
256 y claves derivadas mediante PBKDF2. Publicada en npm.

\entry{picopass}{\href{https://github.com/antruc/picopass}{github.com/antruc/picopass}}

Dispositivo USB HID sobre Raspberry Pi Pico (RP2040) que escribe una cadena
almacenada al conectarse. Firmware en CircuitPython.

%==================================================
% SKILLS
%==================================================

\section{Conocimientos}

% style=unboxed keeps the label inline and lets long lists wrap back to the left
% margin instead of hanging under the label.
\begin{description}[leftmargin=0pt, labelsep=0.5em, itemsep=4pt, topsep=0pt,
                    parsep=0pt, style=unboxed]

    \item[Lenguajes:] Python, TypeScript, JavaScript, HTML, CSS, Bash, VBA

    \item[Frameworks:] React, Vue.js, Svelte, Node.js, Capacitor, Toga y
    Briefcase

    \item[Sistemas operativos:] GNU/Linux, Windows, macOS

    \item[Herramientas:] Git, Visual Studio Code, Android Studio, Arduino,
    Raspberry Pi Pico

    \item[IA:] Ultralytics YOLO, ONNX Runtime, llama.cpp

    \item[Diseño y edición:] GIMP, Inkscape, Kdenlive, Audacity, Canva

    \item[Administración de TI:] Microsoft 365, SharePoint

    \item[SIG:] ArcGIS, QGIS

    \item[Idiomas:] Español (nativo), Inglés (avanzado)

\end{description}

\end{document}
